\documentclass[11pt,a4paper]{article}
\usepackage[utf8]{inputenc}
\usepackage[T1]{fontenc}
\usepackage{amsmath,amssymb}
\usepackage{geometry}
\geometry{a4paper, margin=2.5cm}
\usepackage{xcolor}
\usepackage{booktabs}
\usepackage{array}
\usepackage{enumitem}
\usepackage{fancyhdr}
\usepackage{tcolorbox}
\usepackage{hyperref}

% Color palette
\definecolor{primary}{RGB}{24, 43, 73}      % Dark Navy
\definecolor{secondary}{RGB}{180, 83, 9}    % Amber / Warm Accent
\definecolor{accent}{RGB}{37, 99, 235}      % Blue Accent
\definecolor{bglight}{RGB}{248, 250, 252}   % Light slate background

% Hyperref setup
\hypersetup{
    colorlinks=true,
    linkcolor=primary,
    citecolor=secondary,
    urlcolor=accent,
    pdftitle={Analyse Mecanistique de covtl-pipeline : Nature de Lambda et Trajectoires C1C2},
    pdfauthor={Gemini Notebook}
}

% Tcolorbox setup
\tcbuselibrary{skins,breakable}
\newtcolorbox{summarybox}[1][]{
    colback=bglight,
    colframe=primary,
    coltitle=white,
    fonttitle=\bfseries,
    title=#1,
    arc=3pt,
    boxrule=1pt,
    left=10pt, right=10pt, top=8pt, bottom=8pt,
    breakable
}

\newtcolorbox{keyinsight}[1][]{
    colback=bglight,
    colframe=secondary,
    fonttitle=\bfseries,
    coltitle=white,
    title=#1,
    arc=3pt,
    boxrule=1pt,
    left=10pt, right=10pt, top=8pt, bottom=8pt,
    breakable
}

% Header and Footer
\pagestyle{fancy}
\fancyhf{}
\lhead{\small \textcolor{gray}{\textbf{Modélisation Articulatoire \& Reconnaissance Vocale}}}
\rhead{\small \textcolor{gray}{\textbf{Analyse du modèle \texttt{covtl-pipeline}}}}
\cfoot{\thepage}
\renewcommand{\headrulewidth}{0.4pt}
\renewcommand{\footrulewidth}{0.4pt}

% Title styling
\title{\vspace{-1cm}
\Huge \textbf{\textcolor{primary}{Modélisation Articulatoire Continue et Contrôle Causal dans \texttt{covtl-pipeline}}}\\
\Large \vspace{0.3cm} \textcolor{secondary}{Nature du Paramètre $\lambda$, Extension de la Branche CC et Invariance de $C_2$ dans les Séquences $C_1 C_2$}
}
\author{\textbf{Synthèse Technico-Théorique} \\ Reconnaissance Vocale Neuronale \& Synthèse Articulatoire}
\date{\today}

\begin{document}

\maketitle

\begin{abstract}
\noindent Ce rapport formalise la mécanique interne du module de synthèse articulatoire \texttt{covtl-pipeline} (pilotant \textit{VocalTractLab}) appliqué à l'étude des représentations dans les réseaux de neurones vocaux (notamment Wav2Vec2). Il répond de manière rigoureuse à deux questions fondamentales :
(1) \textbf{La nature du paramètre $\lambda$} : $\lambda$ ne constitue pas une simple pondération ou un fondu enchaîné (\textit{cross-fade}) entre deux états physiques discrets, mais correspond à un déplacement vectoriel continu dans l'espace de contrôle polaire $(\rho, \theta)$ de la branche consonantique (CC), couplé à une extension du sélecteur d'organes à l'union $S_{\text{union}} = S_n \cup S_m$.
(2) \textbf{La dynamique des séquences $C_1 C_2$} : Dans une chaîne consonantique $C_1 C_2$, la consonne contextuelle $C_2$ demeure strictement invariable en position $(\rho_{C2}, \theta_{C2})$ et en timing, tandis que seul le nœud de commande de la consonne cible $C_1$ subit une dérive causale continue. Cette asymétrie garantit une isolation causale parfaite du signal d'entrée.
\end{abstract}

\vspace{0.3cm}

\section{Introduction et Cadre Théorique}

L'analyse du traitement de l'assimilation de lieu d'articulation (ex. \textit{clean pan} prononcé «~clea[m] pan~») par les modèles de reconnaissance automatique de la parole (ASR) repose sur la confrontation entre la parole humaine naturelle et la parole synthétisée. Alors que la parole humaine naturelle comporte des variables confondues (micro-traces coarticulatoires, prosodie, biais de fréquence de mots), la synthèse articulatoire via \texttt{covtl-pipeline} offre un \textbf{instrument d'intervention causale directe}.

Afin de modéliser le passage graduel d'une consonne canonique alvéolaire $/n/$ à une consonne assimilée bilabiale $[m]$, les expérimentations s'appuient sur un continuum contrôlé par un paramètre $\lambda \in [0, 1]$. L'objet de ce document est de formaliser la réalité physique et informatique sous-jacente à ce paramètre $\lambda$ et d'analyser le comportement des séquences consonantiques $C_1 C_2$.

\section{Nature du Paramètre $\lambda$ : Espace Polaire $(\rho, \theta)$ et Union des Sélecteurs}

\subsection{Dépassement du modèle de pondération à deux états}
Une conception naïve de la synthèse d'assimilation consisterait à envisager $\lambda$ comme un simple «~fondu enchaîné~» (\textit{cross-fade}) acoustique ou un mélange pondéré de deux signaux audio préexistants. Cette interprétation est inexacte. Le modèle \texttt{covtl-pipeline} n'effectue aucun mélange de signaux, mais procède à une \textbf{ré-allocation dynamique dans l'espace de contrôle moteur du conduit vocal}.

\subsection{Représentation dans le plan polaire de la branche CC}
Dans le modèle de contrôle articulatoire \texttt{covtl}, la branche consonantique (CC) d'une \textbf{consonne double tautosyllabique} définit chaque cible comme une position complexe $z_c = \rho e^{i\theta}$ dans un plan polaire de coordination. Pour la paire de référence \textit{clean pan} écrite comme la consonne double \texttt{klinpan} (k-l-i-n-p-a-n ; forme sans point --- un point syllabique est une scission C.C, une structure de geste différente, hors du périmètre de la variation $\lambda$) :
\begin{itemize}[leftmargin=*]
    \item \textbf{Cible canonique $/n/$ (coda nasale de la consonne double)} : $z_c^{[n]} = (\rho = 1{,}15,\, \theta = 270^\circ)$ --- la règle de la consonne double du moteur donne à la nasale en cluster la cible alvéolaire de type occlusive.
    \item \textbf{Cible de contexte $C_2 = /p/$ (attaque de la consonne double, fixe)} : $z_c^{[p]} = (\rho = 1{,}06,\, \theta = 50^\circ)$.
\end{itemize}

Le paramètre $\lambda$ déplace la \textbf{coordonnée $\theta$ de la cible $C_1$ vers le $\theta$ de la cible $C_2$} selon le plus court chemin angulaire ($\Delta\theta = +140^\circ$ ici), le canal $\rho$ étant optionnel :
\begin{equation}
\theta_{C1}(\lambda) = \theta_{C1} + \lambda \,\operatorname{wrap}_{(-\pi,\pi]}\!\left(\theta_{C2} - \theta_{C1}\right),
\qquad \rho_{C1}(\lambda) = (1-\lambda)\,\rho_{C1} + \lambda\,\rho_{C2} \;(\text{optionnel})
\end{equation}

\noindent À $\lambda = 0$, la cible de la consonne double est strictement $C_1C_2$ ; à $\lambda = 1$, la cible polaire de $C_1$ a atteint les coordonnées de la cible $C_2$ ($\theta$ exactement ; $\rho$ aussi si activé) --- $C_1'C_2$ : le geste de $C_1$ \emph{pointe désormais vers} la cible de $C_2$, tandis que $C_2$ elle-même ne bouge jamais. À $\lambda = 0{,}5$ avec le canal par défaut ($\theta$ seul), le nœud de commande se situe à $(\rho = 1{,}15,\, \theta = 340^\circ)$, à mi-chemin du trajet angulaire.

\subsection{Union des sélecteurs de la consonne double ($S_{\text{union}}$)}
La machinerie de la consonne double du moteur pilote, sur chaque sous-arc du cluster, \textbf{l'union des sélecteurs des consonnes} :
\begin{itemize}[leftmargin=*]
    \item Sélecteur lingual de $/n/$ : $S_n = \{\text{JA, TCX, TCY, TBX, TBY, TTX, TTY}\}$.
    \item Sélecteur labial de $/p/$ : $S_p = \{\text{JA, LD, LP}\}$.
    \item \textbf{Sélecteur de la consonne double} : $S_{\text{union}} = S_n \cup S_p = \{\text{JA, TCX, TCY, TBX, TBY, TTX, TTY, LD, LP}\}$.
\end{itemize}

Cette machinerie appartient au moteur et est \emph{indépendante de $\lambda$} : la variation ne fait que déplacer les coordonnées de cible de $C_1$ à l'intérieur ($C_1$ conserve son identité nasale --- sélecteur et enveloppe de nasalité --- tandis que sa cible pivote vers le lieu de $C_2$). Pour chaque valeur de $\lambda$, la trajectoire est reconstruite et l'intégralité des 19 paramètres anatomiques du conduit vocal est recalculée trame par trame (à une fréquence de $400\text{ Hz}$).

\begin{summarybox}[Synthèse : Ce que représente physiquement $\lambda$]
Le paramètre $\lambda$ fait pivoter la cible polaire de la nasale $C_1$ vers le lieu de $C_2$ à l'intérieur de la machinerie fixe de la consonne double ($S_{\text{union}}$). À $\lambda = 0{,}5$ (paire de référence \texttt{klinpan}), le système ne superpositionne pas deux sons ; il synthétise un geste hybride unique --- une approximante nasalisée dont la cible se situe à mi-chemin du trajet angulaire ($\theta = 340^\circ$) : le corps de la langue pivote déjà vers le labial (mesuré au milieu de la fenêtre $C_1$ : mâchoire $-2{,}29^\circ$ contre $-2{,}32^\circ$ en canonique, pointe de langue en $x$ $4{,}23$ cm contre $4{,}78$ cm) tandis que le voile du palais reste grand ouvert.
\end{summarybox}

\section{Séquences $C_1 C_2$ : Invariance de $C_2$ et Glissement Causal de $C_1$}

\subsection{Formalisation des paires de consonnes}
Lorsqu'on étudie une séquence consonantique $C_1 C_2$ dans un contexte d'assimilation régressive (\textit{clean pan} où $C_1 = \text{/n/}$ et $C_2 = \text{/p/}$, ou \textit{clean tan} où $C_1 = \text{/n/}$ et $C_2 = \text{/t/}$, réalisés sur le moteur par les clusters tautosyllabiques \textit{klinpan} / \textit{klintan}), la trajectoire de commande de la branche CC est construite selon une asymétrie stricte.

\subsection{Invariance de la consonne contextuelle $C_2$}
La consonne $C_2$ sert de déclencheur phonologique ou de contexte de contrôle. Dans le pipeline \texttt{covtl}, \textbf{la position géométrique $(\rho_{C2}, \theta_{C2})$ et le timing de $C_2$ sont strictement invariants} (vérifié bit à bit à chaque $\lambda$ par le contrôle \texttt{outside\_c1\_arcs\_bitwise\_canon}) :
\begin{itemize}[leftmargin=*]
    \item Pour \textit{clean pan} ($C_2 = \text{/p/}$) : Cible bilabiale fixée à $(\rho = 1{,}06,\, \theta = 50^\circ)$.
    \item Pour \textit{clean tan} ($C_2 = \text{/t/}$) : Cible alvéolaire fixée à $(\rho = 1{,}20,\, \theta = -75^\circ)$ (règle de la consonne double du moteur pour \texttt{klintan}).
\end{itemize}

\subsection{Dérive causale exclusive de $C_1$}
Seule la cible de la consonne nasale $C_1$ varie en fonction de $\lambda$ :
\begin{equation}
\theta_{C1}^{(\lambda)} = \theta_{[n]} + \lambda \,\operatorname{wrap}\!\left(\theta_{C2} - \theta_{[n]}\right)
\end{equation}

\begin{table}[h]
\centering
\caption{Évolution des paramètres de $C_1$ et $C_2$ le long du continuum $\lambda$ (consonne double de référence \texttt{klinpan} ; canal $\theta$ seul, par défaut)}
\vspace{0.2cm}
\small
\begin{tabular}{lcccc}
\toprule
\textbf{Régime $\lambda$} & \textbf{Cible $C_1$ dans $(\rho, \theta)$} & \textbf{Position $C_2$ (\textit{pan})} & \textbf{Nature du Geste $C_1$} & \textbf{Régime d'Assimilation} \\
\midrule
$\lambda = 0{,}00$ & $(1{,}15,\, 270^\circ)$ & $(1{,}06,\, 50^\circ)$ & Nasale alvéolaire de type occlusive & Forme canonique /n/ \\
$\lambda = 0{,}50$ & $(1{,}15,\, 340^\circ)$ & $(1{,}06,\, 50^\circ)$ & Approximante nasalisée & Geste hybride / ambigu \\
$\lambda = 0{,}75$ & $(1{,}15,\, 15^\circ)$ & $(1{,}06,\, 50^\circ)$ & Nasale quasi-labiale & Assimilation partielle \\
$\lambda = 1{,}00$ & $(1{,}15,\, 50^\circ)$ & $(1{,}06,\, 50^\circ)$ & Fermeture nasale pointant vers le labial & Partage de cible ($\theta$ ; $\rho$ aussi avec le canal rho) \\
\bottomrule
\end{tabular}
\end{table}

\subsection{Portée pour la rigueur du contrôle expérimental}
Cette asymétrie garantit que la modification physique appliquée sur le segment $C_1$ est \textbf{strictement identique sur le plan acoustique et articulatoire}, que la consonne suivante $C_2$ soit un $/p/$ (contexte viable) ou un $/t/$ (contexte non viable).

\begin{keyinsight}[Implication Causale Majeure]
Par construction, le signal audio modifiant $C_1$ ne contient aucune information sur le contexte $C_2$. Si le réseau de neurones (Wav2Vec2) réagissait différemment entre contextes viables et non viables sur ce signal synthétique pur, cela prouverait sans ambiguïté l'existence d'une règle de licence phonologique contextuelle au sein du réseau.
\end{keyinsight}

\section{Implications pour la Modélisation Cognitives et les Réseaux de Neurones}

Les simulations conduites sur ce continuum synthétique pour différentes valeurs de $\lambda$ ($\lambda \in \{1{,}0\,;\, 0{,}75\,;\, 0{,}625\}$) apportent un résultat fondamental pour l'interprétation du modèle Wav2Vec2 :

\begin{enumerate}[leftmargin=*]
    \item \textbf{Absence d'asymétrie de viabilité} : À chaque valeur de $\lambda$, le taux de compensation comportemental et la probabilité interne de la forme sous-jacente $P(n)$ restent \textbf{strictement équivalents} en contexte viable (\textit{clean pan}) et en contexte non viable (\textit{clean tan}).
    \item \textbf{Démonstration de la nature de l'encodeur} : L'encodeur de Wav2Vec2 (couches CNN à T6) agit comme un \textbf{détecteur acoustique ascendant (\textit{bottom-up})}. Il ne possède aucune règle de grammaire phonologique abstraite autonome.
    \item \textbf{Explication de l'effet sur parole humaine} : La compensation asymétrique mesurée sur la parole humaine dans l'étude de Pouw et al. (2024) provient uniquement de la présence de \textbf{micro-traces physiques coarticulatoires} dans la voix humaine et des \textbf{statistiques orthographiques du décodeur CTC} (couches T10 à T12).
\end{enumerate}

\section{Conclusion}

La formalisation du module \texttt{covtl-pipeline} démontre la grande rigueur expérimentale offerte par la synthèse articulatoire :
1. Le paramètre $\lambda$ est une \textbf{rotation continue de la cible polaire de $C_1$ vers la cible de $C_2$} dans l'espace polaire $(\rho, \theta)$ de la consonne double tautosyllabique ($\theta$ toujours, $\rho$ optionnellement), à l'intérieur de la machinerie fixe du cluster (union des sélecteurs $S_{\text{union}}$) ; à $\lambda = 1$, la cible de $C_1$ égale les coordonnées de la cible $C_2$. Chaque $\lambda$ est une véritable consonne double re-synthétisée (le mélange paramétrique trame à trame de l'expérience historique du kit reste disponible comme approximation au niveau moteur).
2. Dans une séquence $C_1 C_2$, la manipulation est \textbf{strictement localisée sur $C_1$}, maintenant $C_2$ et tout le contexte invariants bit à bit.

Ce cadre expérimental permet de dissocier définitivement la sensibilité acoustique directe de la correction orthographique supervisée dans les architectures de traitement automatique de la parole.

\end{document}
